\documentclass[11pt]{article}
\usepackage[margin=1in]{geometry}
\usepackage{amsmath,amssymb,amsthm}
\usepackage[T1]{fontenc}
\usepackage{lmodern}
\usepackage{graphicx}
\usepackage[hidelinks]{hyperref}
\usepackage{xcolor}

\newtheorem{theorem}{Theorem}
\newtheorem{corollary}[theorem]{Corollary}
\newcommand{\N}{\mathbb N}
\newcommand{\sid}{\operatorname{sid}}

\title{A Density-Deficit Obstruction for\\
Near-Full Sidon Subsystems}
\author{Candidate partial result for arXiv:1504.05290}
\date{9 August 2026}

\begin{document}
\maketitle

\begin{center}
\fcolorbox{black}{gray!8}{\parbox{0.88\linewidth}{
\textbf{Status: candidate partial result, likely valid.}
Problem 26 remains open for each fixed positive density deficit.}}
\end{center}

\section{Source problem}

Bourgain and Lewko ask the following in Problem 26 of
\cite{BL} (arXiv PDF page 20).

\begin{quote}
Does there exist a constant $\gamma=\gamma(M,C,\varepsilon)$ such that for
any orthonormal system of size $n$, uniformly bounded by $M$ and satisfying
the $\psi_2(C)$ condition, there exists $A\subseteq[n]$ with
$|A|\geq(1-\varepsilon)n$ such that
\[
 \left\|\sum_{j\in A}a_j\phi_j\right\|_{L^\infty}
 \geq \gamma\sum_{j\in A}|a_j| ?
\]
\end{quote}

\begin{figure}[h]
  \centering
  \includegraphics[width=\textwidth]{figures/open_problem_crop.png}
  \caption{Problem 26 in arXiv:1504.05290, PDF page 20.}
\end{figure}

For a finite system $\mathcal A=(f_j)_{j\in A}$, write
\[
 \sid(\mathcal A)=
 \inf_{(a_j)\ne0}
 \frac{\|\sum_{j\in A}a_jf_j\|_\infty}{\sum_{j\in A}|a_j|}
\]
for its optimal Sidon lower constant.

\section{Proof intuition}

The Bourgain--Lewko counterexample has two coupled pieces.  Its Rademacher
piece would normally make equal coefficients large, but all functions share
a factor which damps a Rademacher configuration according to the square of
the total sign sum.  If $k$ coordinates have been deleted, those missing
coordinates can neutralize at most $k$ units of the retained sign sum.  Past
that threshold the common factor caps the retained contribution by
$n/\sqrt{\log n}$.  The Walsh--Rudin--Shapiro perturbation is also small:
the full signed Walsh sum is $O(\sqrt n)$, while deleting $k$ summands costs
at most $k$.

Equal coefficients therefore witness an upper Sidon constant of
$k/(n-k)+o(1)$.  Hence a density-$1-\varepsilon$ subsystem can have no
universal constant larger than $\varepsilon/(1-\varepsilon)$, and deleting
only $o(n)$ terms never suffices.  This does not decide whether a positive
constant of order at most $\varepsilon$ always exists for fixed
$\varepsilon$.

\section{The quantitative obstruction}

We use the real orthonormal systems from Section 9 of \cite{BL}.  Fix large
$n$, put $L=\log n$ and $\delta=\sqrt{L/n}$, let $r_1,\ldots,r_n$ be
independent Rademacher functions, and choose Walsh functions $W_i$ and signs
$\sigma_i$ so that
\begin{equation}\label{eq:RS}
 \left\|\sum_{i=1}^n\sigma_iW_i\right\|_\infty\leq6\sqrt n.
\end{equation}
The probability density used in \cite{BL} is
\[
 \Psi=\left(1+\frac{L}{n}\right)^{-1}
 \left(1+\frac{L}{n^2}\left(\sum_{i=1}^nr_i\right)^2\right).
\]
For $1\leq i\leq n$ their functions can consequently be written exactly as
\begin{equation}\label{eq:phi}
 \phi_i=D(r)(r_i-\delta\sigma_iW_i),\qquad
 D(r)=\left(1+\frac{L}{n^2}\left(\sum_{i=1}^nr_i\right)^2\right)^{-1/2}.
\end{equation}
Together with the additional function $\phi_0$, these form an orthonormal
system $\Phi^{(n)}=\{\phi_0,\phi_1,\ldots,\phi_n\}$, uniformly bounded by
$7$ and satisfying $\psi_2(C_0)$ for a universal $C_0$ \cite{BL}.

\begin{theorem}\label{thm:quant}
Let $\mathcal A\subseteq\Phi^{(n)}$, put
\[
 B=\{i\in\{1,\ldots,n\}:\phi_i\in\mathcal A\},\qquad
 m=|B|,\qquad k=n-m,
\]
and suppose $m>0$.  Then
\begin{equation}\label{eq:bound}
 \sid(\mathcal A)
 \leq
 \frac{k+n/\sqrt{\log n}+6\sqrt{\log n}
       +k\sqrt{\log n/n}}{m}.
\end{equation}
In particular, for every fixed $0<\varepsilon<1$ and every choice
$\mathcal A_n\subseteq\Phi^{(n)}$ satisfying
$|\mathcal A_n|\geq(1-\varepsilon)(n+1)$,
\begin{equation}\label{eq:limsup}
 \limsup_{n\to\infty}\sid(\mathcal A_n)
 \leq\frac{\varepsilon}{1-\varepsilon}.
\end{equation}
\end{theorem}

\begin{proof}
Assign coefficient $1$ to the functions $\phi_i$ with $i\in B$ and
coefficient $0$ to $\phi_0$ if it belongs to $\mathcal A$.  This coefficient
vector has $\ell_1$ norm $m$.  Write $K=\{1,\ldots,n\}\setminus B$ and
\[
 S_B=\sum_{i\in B}r_i,\qquad S_K=\sum_{i\in K}r_i,
 \qquad P_B=\sum_{i\in B}\sigma_iW_i.
\]
Since $|S_K|\leq k$,
\begin{equation}\label{eq:total}
 |S_B+S_K|\geq (|S_B|-k)_+.
\end{equation}
Moreover, \eqref{eq:RS} and $|W_i|=1$ give
\begin{equation}\label{eq:PB}
 |P_B|
 \leq\left|\sum_{i=1}^n\sigma_iW_i\right|
      +\left|\sum_{i\in K}\sigma_iW_i\right|
 \leq6\sqrt n+k.
\end{equation}
Equations \eqref{eq:phi}--\eqref{eq:PB}, with $x=|S_B|$, imply pointwise
\begin{align*}
 \left|\sum_{i\in B}\phi_i\right|
 &\leq
 \frac{x}{\sqrt{1+(L/n^2)(x-k)_+^2}}
 +\delta(6\sqrt n+k).
\end{align*}
The first term is at most $k+n/\sqrt L$.  Indeed, it is at most $k$ when
$x\leq k$; when $x=k+y>k$, it is bounded by
\[
 \frac{k+y}{\sqrt{1+(L/n^2)y^2}}
 \leq k+\frac{y}{\sqrt{1+(L/n^2)y^2}}
 \leq k+\frac n{\sqrt L}.
\]
Taking the supremum and using
$\delta(6\sqrt n+k)=6\sqrt L+k\sqrt{L/n}$ proves \eqref{eq:bound}.

If $|\mathcal A_n|\geq(1-\varepsilon)(n+1)$, then
\[
 m\geq |\mathcal A_n|-1
 \geq(1-\varepsilon)n-\varepsilon,
 \qquad
 k\leq\varepsilon n+\varepsilon.
\]
Divide \eqref{eq:bound} by $m$ and let $n\to\infty$.  Every term except
$k/m$ tends to zero, while
$\limsup k/m\leq\varepsilon/(1-\varepsilon)$.  This proves
\eqref{eq:limsup}.
\end{proof}

\begin{corollary}
There are universally bounded orthonormal $\psi_2(C_0)$ systems for which no
subsystem obtained by deleting $o(n)$ functions has a Sidon lower constant
bounded away from zero.  Consequently, if Problem 26 has an affirmative
answer, its best possible universal dependence for $M=7,C=C_0$ satisfies
\[
 \gamma(7,C_0,\varepsilon)leq
 \frac{\varepsilon}{1-\varepsilon}
 \qquad(0<\varepsilon<1).
\]
\end{corollary}

\begin{proof}
For the first assertion apply Theorem~\ref{thm:quant} with $k=o(n)$.
For the second, a constant asserted by Problem 26 must work for every member
of the Bourgain--Lewko family.  The bound \eqref{eq:limsup} applies to every
subsystem of the required size, so no larger universal constant is possible.
\end{proof}

\section{Limitations and novelty check}

The theorem is not a counterexample to Problem 26 for fixed
$\varepsilon>0$: the right side of \eqref{eq:limsup} is then positive, and
the problem explicitly allows $\gamma$ to depend on $\varepsilon$.  What is
settled is the stronger density-one version with $o(n)$ deletions and an
explicit linear upper restriction on any possible fixed-$\varepsilon$ answer.

A bounded arXiv and local full-source search used arXiv:1504.05290, papers
citing it, the exact phrase of Problem 26, and the terms ``density-one Sidon
subsystem'', ``Sidon lower constant'', and ``uniformly bounded $\psi_2$
orthonormal system''.  Pisier's later work resolves the two-fold tensor
question but not Problem 26.  No source stating Theorem~\ref{thm:quant} or
its corollary was located.  This supports, but does not certify, novelty.

\section{Verification and human-review focus}

The proof is elementary once the source construction is written as
\eqref{eq:phi}.  The common-factor identity, the two triangle inequalities
\eqref{eq:total}--\eqref{eq:PB}, and all limiting quantifiers were checked
separately in the accompanying verification note.  No numerical computation
is used.

Human review should focus on the exact rewriting of the density $\Psi$ and on
the use of zero coefficient for $\phi_0$ when the tested subsystem contains
that function.

\begin{thebibliography}{9}
\bibitem{BL}
J. Bourgain and M. Lewko,
\emph{Sidonicity and variants of Kaczmarz's problem},
Ann. Inst. Fourier (Grenoble) 67 (2017), no.~3, 1321--1352;
arXiv:1504.05290.
\end{thebibliography}

\end{document}
